\chapter{Future Scope}\label{ch:future}

\begin{enumerate}
\item \textbf{A like-for-like comparison with WaW.} Run the pipeline on the PhenoBench validation
  set~\cite{weyler2024phenobench} (pixel labels) and report weed IoU, the metric of Chong et
  al.~\cite{chong2025zeroshot}.
\item \textbf{Species-aware features instead of text labels.} Replace the detector's words with
  BioCLIP~\cite{stevens2024bioclip} features of each plant, as WaW does, and keep the safety gate on
  top.
\item \textbf{A ``safe mode'' for rule 4.} Limit the off-row rule by plant size (Table~\ref{tab:f1}),
  tested on new held-out data, so that a missed crop box can no longer cause a whole crop plant to be sprayed.
\item \textbf{Few-shot adaptation.} Adapt the detector with a handful of labelled field photos, as
  Singh et al.~\cite{singh2025fewshot} suggest, and measure how many labels close the gap to supervised
  detectors~\cite{rahman2023cotton,dang2023yoloweeds}.
\item \textbf{Using the row over time.} Exploit image sequences along the row~\cite{lottes2018fcn,lottes2020stem}
  and automatically labelled inter-row weeds~\cite{bah2018unsupervised} to calibrate the rules per field.
\item \textbf{Farmer-facing agent and interface.} A conversational assistant over the same verified
  facts, in Indian languages, on a mobile-first web interface.
\item \textbf{Real deployment.} GPS from drone or sprayer metadata, nozzle latency, a field trial with
  an agronomist, and the other applications of the original proposal (disease, yield, canopy cover,
  nutrient stress).
\end{enumerate}
